% !TEX root = ../main.tex
\addcontentsline{toc}{chapter}{Conclusion}
\chapter*{Conclusion}

This thesis started as a problem proposed by one of my supervisors, Dr. Joost J. Joosten, that consisted in generating 
terms of specific types in simply typed combinatory logic $\textbf{CL}_\to$.
Our initial work building a language model to generate those terms made the limitations of this models very clear.
First, linear generation forces the model to overcommit too soon to certain solution paths, resulting in unsatisfactory results, this was noted by researchers in Logical Intelligence in \cite{logicalintelligence}.
Second, that syntax does not hold all the information the model needed to generate the desired terms.

Our study of Lafforgue's framework \cite{lafforgue2022, lafforgue2024} 
connecting Topos Theory \cite{lawvere1971} with Artificial 
Intelligence, lead us to consider working with categorical representations 
of our type theory in order to leverage its geometrical properties to 
give the model a complete picture of the semantics of the theory.

The first representation we found promising Cartmell's Generalized Algebraic
Theories and Contextual Categories \cite{cartmell1978, cartmell1986}, whose term models
sit close enough to the syntax of a theory to be
computationally manageable, while still carrying the 
geometric structure needed for a GNN-based approach. 
In Section \ref{sect:relevance-of-the-context-category} we revised the definition of model of a GAT, that originally is thought as a mapping to the category of sets, families of sets, families of families of sets, \dots;
but we thought it was more natural to think of an interpretation as a mapping directly to a contextual category.
With this revised notion of model, we proposed a new proof for the initiality conjecture for GATs.
Then we designed the Expansion Algorithm to construct sections within the term model, equivalent to generating terms of a given type.
This model was thought to work together with a GNN in order to surpass the 
complexity barrier.

To handle a more expressive type theory, we turned to Streicher's 
work on the Calculus of Constructions and Doctrines of Constructions \cite{streicher}, 
which provided both a proof of the initiality conjecture for CoC 
and a framework for extending type theories beyond GATs. 
In Section \ref{sect:extending-calculus-of-constructions-and-its-interpretation}, we extended this framework 
to CoC enriched with a GAT signature, constructed the corresponding 
term model, following ideas 
from \cite{bezem2021}, and proposed an outline of the proof of its initiality based on the original proof.
We then extended the Expansion Algorithm to this 
richer setting. To do this, we extended the concept of relativization to Doctrines of Constructions. 
In the course of this work, we also identified a 
connection between the concept of collections of types in Doctrines 
of Constructions and the topos-theoretic notion of subobject 
classifier \cite{maclane1992}, which we leave as a direction for further 
investigation.

Several directions remain open. On the theoretical side, 
it remains to formally extend the term model construction
and prove the initiality conjecture for Doctrines of 
Constructions with $\Sigma$-types, identity types, and a universe 
hierarchy, as described in \cite{streicher}, and ultimately for the full 
Calculus of Inductive Constructions \cite{coquand1990} \cite{paulin-mohring1993} and the Lean type 
theory \cite{demoura2015}. The initiality conjecture for full Martin-L\"of type theory 
has been formalized in Agda by Brunerie and Lumsdaine \cite{bruneire2020}, 
and a similar formalization effort for our extended setting 
would be a natural next step.

On the practical side, the full implementation of the GNN-based 
Conjecturer-Prover model for guiding the Expansion Algorithm 
remains the primary target for future work. We believe this 
architecture, by combining the geometric structure of categorical 
term models with the pruning power of graph neural networks, 
offers a principled and promising alternative to purely syntactic 
approaches to autonomous theorem proving.